% 第一章 投资概要与核心观点

\chapter{投资概要与核心观点}

\section{投资结论速览：超配/标配/关注/回避四类评级}

本报告基于“业绩验证—产业映射—政策加持—估值持仓”四维度交叉验证框架，对A股31个一级行业、120+细分板块进行系统性扫描，识别景气度真实但市场预期尚未充分定价的投资机会。综合景气度兑现度、预期差空间、估值安全边际三大核心指标，我们将覆盖板块划分为\textbf{超配、标配、关注、回避}四类评级。

\textbf{超配板块（3个）：功率半导体、电力公用事业、智驾链零部件}

三大板块共同特征是景气度已通过业绩端验证、产业趋势明确、政策面持续加持，而市场认知仍停留在传统框架下，存在显著预期差。功率半导体受益于AI算力+能源转型双轮驱动，2026年Q1行业营收同比增长28\%，但板块持仓集中度仅处于历史35\%分位；电力公用事业在电价机制改革与AI用电增量双重利好下，盈利稳定性大幅提升，当前PE仅12--15倍，处于历史估值底部区域；智驾链零部件迎来L3/L4渗透率拐点，2026年高阶智驾渗透率有望从15\%跃升至35\%，而产业链多数标的仍按传统汽车零部件估值定价。

\textbf{标配板块（2个）：半导体设备材料、医疗器械}

半导体设备材料受益于国内晶圆厂扩产周期与国产替代加速，2026年设备招标量同比增长40\%以上，但受制于海外制裁不确定性，市场给予一定风险折价；医疗器械板块经过三年调整，估值与持仓均处于历史低位，消费医疗与设备更新周期带来边际改善，但整体弹性弱于超配板块。

\textbf{关注板块（2个）：动力电池龙头、创新药龙头}

两大板块均处于左侧布局窗口，行业基本面已出现边际改善信号，但趋势性拐点尚未确认。动力电池龙头价格战接近尾声，盈利能力触底回升，海外市场打开第二增长曲线；创新药龙头受益于医保谈判边际缓和与出海加速，但研发管线兑现度仍需观察。

\textbf{回避板块（3个）：6G、低空经济、卫星互联网}

三大主题板块景气度尚未通过业绩端验证，政策预期透支严重，商业模式模糊，属于典型的“博弈型”甚至“陷阱型”板块。当前估值中已包含大量远期乐观假设，而产业落地进度普遍慢于市场预期，建议低配或回避。

\begin{exhibitbox}[图表1-1：板块评级全景图——超配/标配/关注/回避四象限]
\centering
\vspace{0.3cm}
\begin{tikzpicture}[scale=0.95]
% 坐标轴
\draw[->, thick, navy] (0,0) -- (9,0) node[right] {\small 景气度兑现度};
\draw[->, thick, navy] (0,0) -- (0,7) node[above] {\small 预期差空间};

% 坐标轴刻度
\foreach \x/\xlabel in {0/0\%, 2.25/25\%, 4.5/50\%, 6.75/75\%, 9/100\%}
  \draw[thick, navy] (\x, 0.1) -- (\x, -0.1) node[below, font=\scriptsize, color=steelgrey] {\xlabel};
\foreach \y/\ylabel in {0/0\%, 1.75/25\%, 3.5/50\%, 5.25/75\%, 7/100\%}
  \draw[thick, navy] (-0.1, \y) -- (0.1, \y) node[left, font=\scriptsize, color=steelgrey] {\ylabel};

% 四象限分割线
\draw[dashed, rulegrey] (4.5,0) -- (4.5,7);
\draw[dashed, rulegrey] (0,3.5) -- (9,3.5);

% 象限标签
\node[font=\bfseries\color{navy}, align=center] at (2.25, 6.3) {I 象限\\超配};
\node[font=\bfseries\color{steelgrey}, align=center] at (6.75, 6.3) {II 象限\\标配};
\node[font=\bfseries\color{riskamber}, align=center] at (2.25, 1.2) {III 象限\\关注};
\node[font=\bfseries\color{riskred}, align=center] at (6.75, 1.2) {IV 象限\\回避};

% 超配板块（标签分散角度，避免重叠）
\node[circle, fill=navy!50, draw=navy, inner sep=4pt, label={[font=\scriptsize, xshift=-2pt, yshift=2pt]135:功率半导体}] at (1.5, 5.2) {};
\node[circle, fill=navy!50, draw=navy, inner sep=4pt, label={[font=\scriptsize, yshift=-4pt]-90:电力公用事业}] at (2.8, 4.5) {};
\node[circle, fill=navy!50, draw=navy, inner sep=4pt, label={[font=\scriptsize, xshift=2pt, yshift=2pt]45:智驾链零部件}] at (3.8, 5.0) {};

% 标配板块（标签左右分散）
\node[circle, fill=steelgrey!50, draw=steelgrey, inner sep=4pt, label={[font=\scriptsize, xshift=-2pt, yshift=2pt]135:半导体设备材料}] at (5.5, 5.2) {};
\node[circle, fill=steelgrey!50, draw=steelgrey, inner sep=4pt, label={[font=\scriptsize, xshift=2pt, yshift=2pt]45:医疗器械}] at (7.0, 4.5) {};

% 关注板块（标签左下/右下分散）
\node[circle, fill=riskamber!50, draw=riskamber, inner sep=4pt, label={[font=\scriptsize, xshift=-2pt, yshift=2pt]135:动力电池龙头}] at (1.8, 2.8) {};
\node[circle, fill=riskamber!50, draw=riskamber, inner sep=4pt, label={[font=\scriptsize, xshift=2pt, yshift=-2pt]-45:创新药龙头}] at (3.2, 2.2) {};

% 回避板块（标签左/下/右分散）
\node[circle, fill=riskred!50, draw=riskred, inner sep=4pt, label={[font=\scriptsize, xshift=-2pt, yshift=2pt]135:6G}] at (5.5, 2.8) {};
\node[circle, fill=riskred!50, draw=riskred, inner sep=4pt, label={[font=\scriptsize, yshift=-4pt]-90:低空经济}] at (7.0, 2.2) {};
\node[circle, fill=riskred!50, draw=riskred, inner sep=4pt, label={[font=\scriptsize, xshift=2pt, yshift=2pt]45:卫星互联网}] at (8.0, 2.8) {};

% 箭头说明（移至右侧空白区，避免与气泡/标签重叠）
\draw[->, thick, nvgreen] (6.2, 0.6) -- (8.3, 4.2) node[midway, above, sloped, font=\scriptsize, yshift=2pt] {超额收益递增方向};

\end{tikzpicture}
\vspace{0.2cm}
\sourcenote{AStock研究团队绘制。气泡位置为示意性标注，具体数据详见正文各章节。}
\end{exhibitbox}

\section{核心投资逻辑：景气度分层与预期差定价}

本报告的核心投资逻辑建立在\textbf{“景气度分层+预期差定价”}双支柱框架之上，旨在解决传统景气度投资“看对行业、选错时点”的普遍痛点。

\subsection{景气度三维验证体系}

我们摒弃单一维度的景气度判断方式，构建业绩验证、产业映射、政策加持三维交叉验证体系，从根本上规避“伪景气”陷阱。业绩验证关注财报层面的收入与利润增速、毛利率变化、订单能见度等硬指标；产业映射追踪产业链上下游的产能利用率、价格走势、库存水平等领先指标；政策加持评估产业政策的力度、持续性与实际落地效果。只有三维度同时指向景气上行的板块，才纳入我们的核心配置池。

截至2026年Q1，在我们覆盖的120+细分板块中，同时满足三维验证的板块仅13个，占比约10\%，说明真正的高景气度是稀缺资源，这也从侧面解释了为何2026年以来市场呈现极致结构性行情。

\subsection{预期差定价：超额收益的核心来源}

景气度本身不直接带来超额收益，\textbf{预期差才是超额收益的核心来源}。我们引入“拥挤度—兑现度”二维矩阵，将板块划分为四类：

\begin{itemize}[leftmargin=2em]
\item \textbf{明星型}（高兑现+低拥挤）：景气度真实且市场认知不足，是最佳配置方向，典型如功率半导体、电力公用事业
\item \textbf{潜力型}（低兑现+低拥挤）：景气度尚在早期，赔率高但胜率不足，适合左侧布局，典型如创新药龙头
\item \textbf{博弈型}（高兑现+高拥挤）：景气度真实但交易过于拥挤，容易出现阶段性回调，典型如AI算力板块
\item \textbf{陷阱型}（低兑现+高拥挤）：景气度未验证但市场预期极高，是价值陷阱高发区，典型如6G、低空经济
\end{itemize}

\begin{exhibitbox}[图表1-2："拥挤度—兑现度"二维矩阵与板块分类]
\centering
\vspace{0.3cm}
\begin{tikzpicture}[scale=1.05]
% 象限背景填充（置于最底层，避免覆盖坐标轴和分割线）
\fill[riskred!12, draw=none] (0,3.5) rectangle (4.5,7);
\fill[riskblue!12, draw=none] (4.5,3.5) rectangle (9,7);
\fill[riskamber!12, draw=none] (0,0) rectangle (4.5,3.5);
\fill[riskgreen!12, draw=none] (4.5,0) rectangle (9,3.5);

% 坐标轴
\draw[->, thick, navy] (0,0) -- (9,0) node[right] {\small 景气度兑现度};
\draw[->, thick, navy] (0,0) -- (0,7) node[above] {\small 交易拥挤度};

% 坐标轴刻度
\draw[thick, navy] (0,0.12) -- (0,-0.12) node[below, font=\scriptsize, yshift=-2pt] {低};
\draw[thick, navy] (4.5,0.12) -- (4.5,-0.12) node[below, font=\scriptsize, yshift=-2pt] {中};
\draw[thick, navy] (9,0.12) -- (9,-0.12) node[below, font=\scriptsize, yshift=-2pt] {高};
\draw[thick, navy] (-0.12,0) -- (0.12,0) node[left, font=\scriptsize, xshift=-2pt] {低};
\draw[thick, navy] (-0.12,3.5) -- (0.12,3.5) node[left, font=\scriptsize, xshift=-2pt] {中};
\draw[thick, navy] (-0.12,7) -- (0.12,7) node[left, font=\scriptsize, xshift=-2pt] {高};

% 四象限分割线
\draw[dashed, rulegrey] (4.5,0) -- (4.5,7);
\draw[dashed, rulegrey] (0,3.5) -- (9,3.5);

% 象限标签 — 陷阱型（左上：低兑现 + 高拥挤）
\node[font=\bfseries\color{riskred}, align=center] at (2.25, 6.35) {陷阱型};
\node[font=\footnotesize\color{steelgrey!80!black}, align=center] at (2.25, 5.7) {低兑现 + 高拥挤\\价值陷阱高发};

% 象限标签 — 博弈型（右上：高兑现 + 高拥挤）
\node[font=\bfseries\color{riskblue}, align=center] at (6.75, 6.35) {博弈型};
\node[font=\footnotesize\color{steelgrey!80!black}, align=center] at (6.75, 5.7) {高兑现 + 高拥挤\\波段操作策略};

% 象限标签 — 潜力型（左下：低兑现 + 低拥挤）
\node[font=\bfseries\color{riskamber}, align=center] at (2.25, 3.1) {潜力型};
\node[font=\footnotesize\color{steelgrey!80!black}, align=center] at (2.25, 2.65) {低兑现 + 低拥挤\\左侧布局机会};

% 象限标签 — 明星型（右下：高兑现 + 低拥挤）
\node[font=\bfseries\color{riskgreen}, align=center] at (6.75, 2.9) {明星型};
\node[font=\footnotesize\color{steelgrey!80!black}, align=center] at (6.75, 2.2) {高兑现 + 低拥挤\\最优配置方向};

% 陷阱型板块（左上，标签统一朝右）
\node[circle, fill=riskred, inner sep=3pt] at (1.0, 5.0) {};
\node[font=\scriptsize, anchor=west] at (1.5, 5.0) {卫星互联网};
\node[circle, fill=riskred, inner sep=3pt] at (1.5, 4.4) {};
\node[font=\scriptsize, anchor=west] at (2.0, 4.4) {低空经济};
\node[circle, fill=riskred, inner sep=3pt] at (0.8, 3.9) {};
\node[font=\scriptsize, anchor=west] at (1.3, 3.9) {6G};

% 博弈型板块（右上，标签统一朝右）
\node[circle, fill=riskblue, inner sep=4pt] at (5.3, 5.2) {};
\node[font=\scriptsize, anchor=west] at (5.9, 5.2) {AI算力};
\node[circle, fill=riskblue, inner sep=3.5pt] at (5.0, 4.3) {};
\node[font=\scriptsize, anchor=west] at (5.6, 4.3) {智驾链零部件};

% 潜力型板块（左下，标签统一朝右）
\node[circle, fill=riskamber, inner sep=3.5pt] at (0.8, 1.95) {};
\node[font=\scriptsize, anchor=west] at (1.3, 1.95) {创新药龙头};
\node[circle, fill=riskamber, inner sep=4pt] at (0.8, 1.15) {};
\node[font=\scriptsize, anchor=west] at (1.3, 1.15) {动力电池龙头};
\node[circle, fill=riskamber, inner sep=3.5pt] at (2.0, 0.5) {};
\node[font=\scriptsize, anchor=west] at (2.5, 0.5) {医疗器械};

% 明星型板块（右下，标签统一朝右）
\node[circle, fill=riskgreen, inner sep=3.5pt] at (5.3, 1.6) {};
\node[font=\scriptsize, anchor=west] at (5.9, 1.6) {功率半导体};
\node[circle, fill=riskgreen, inner sep=4pt] at (5.3, 0.8) {};
\node[font=\scriptsize, anchor=west] at (5.9, 0.8) {电力公用事业};

\end{tikzpicture}
\vspace{0.2cm}
\sourcenote{AStock研究团队绘制。气泡大小代表板块市值权重；位置为示意性标注，具体数据详见正文各章节。}
\end{exhibitbox}

我们的核心策略是\textbf{超配明星型、关注潜力型、回避陷阱型、博弈型波段操作}。当前市场的核心矛盾在于，大量资金集中在少数热门赛道导致博弈型板块拥挤度极高，而部分景气度真实的明星型板块反而因关注不足存在显著预期差。本报告的核心价值就在于系统性挖掘这些被市场忽视的明星型板块。

\section{推荐组合一览：20只核心标的评分排序}

基于上述研究框架，我们从7大重点覆盖板块中精选20只核心标的，构建“景气度Alpha组合”。选股标准包括：行业地位靠前、业绩确定性高、预期差显著、估值具备安全边际。

我们采用五维评分体系对标的进行综合排序：景气度兑现度（权重30\%）、预期差空间（权重25\%）、产业链卡位（权重20\%）、估值安全边际（权重15\%）、催化确定性（权重10\%）。

\begin{exhibitbox}[图表1-3：20只核心标的评分排序与推荐组合]
\scriptsize
\setlength{\tabcolsep}{3.5pt}
\centering
\vspace{0.3cm}
\begin{tabularx}{\textwidth}{@{}c X l r r r r r @{}}
\toprule
\textbf{排名} & \textbf{标的名称} & \textbf{所属板块} & \textbf{景气度} & \textbf{预期差} & \textbf{产业链} & \textbf{估值} & \textbf{综合} \\
& & & \textbf{(30\%)} & \textbf{(25\%)} & \textbf{卡位(20\%)} & \textbf{安全(15\%)} & \textbf{评分} \\
\midrule
\multicolumn{8}{@{}l}{\textbf{超配组合（7只）}} \\
1 & 时代电气 & 功率半导体 & 9.2 & 9.5 & 9.0 & 8.0 & 9.1 \\
2 & 德赛西威 & 智驾链零部件 & 9.0 & 9.2 & 9.5 & 7.5 & 8.9 \\
3 & 长江电力 & 电力公用事业 & 8.0 & 9.8 & 9.5 & 9.0 & 8.9 \\
4 & 斯达半导 & 功率半导体 & 9.3 & 8.8 & 8.5 & 7.0 & 8.7 \\
5 & 国电南瑞 & 电力公用事业 & 8.2 & 9.0 & 9.2 & 8.5 & 8.7 \\
6 & 伯特利 & 智驾链零部件 & 8.8 & 9.0 & 8.8 & 7.8 & 8.7 \\
7 & 扬杰科技 & 功率半导体 & 8.5 & 8.5 & 8.0 & 8.2 & 8.3 \\
\midrule
\multicolumn{8}{@{}l}{\textbf{标配组合（7只）}} \\
8 & 中微公司 & 半导体设备材料 & 9.5 & 7.5 & 9.2 & 6.5 & 8.4 \\
9 & 迈瑞医疗 & 医疗器械 & 7.8 & 8.5 & 9.0 & 8.5 & 8.3 \\
10 & 北方华创 & 半导体设备材料 & 9.3 & 7.2 & 9.5 & 6.0 & 8.2 \\
11 & 联影医疗 & 医疗器械 & 7.5 & 8.8 & 8.5 & 7.5 & 8.0 \\
12 & 拓荆科技 & 半导体设备材料 & 9.0 & 7.8 & 8.0 & 5.5 & 8.0 \\
13 & 宁德时代 & 动力电池龙头 & 7.0 & 8.2 & 9.5 & 8.8 & 8.0 \\
14 & 恒瑞医药 & 创新药龙头 & 6.5 & 8.5 & 9.0 & 7.0 & 7.5 \\
\midrule
\multicolumn{8}{@{}l}{\textbf{关注组合（6只）}} \\
15 & 华润微 & 功率半导体 & 8.0 & 7.5 & 8.0 & 7.8 & 7.9 \\
16 & 比亚迪 & 动力电池龙头 & 7.2 & 7.5 & 8.8 & 8.0 & 7.8 \\
17 & 东微半导 & 功率半导体 & 8.5 & 7.0 & 7.0 & 7.2 & 7.5 \\
18 & 药明康德 & 创新药龙头 & 6.8 & 8.0 & 8.5 & 7.5 & 7.5 \\
19 & 华大智造 & 医疗器械 & 7.0 & 8.2 & 7.5 & 7.0 & 7.4 \\
20 & 新洁能 & 功率半导体 & 7.8 & 7.0 & 7.2 & 7.5 & 7.4 \\
\bottomrule
\end{tabularx}
\vspace{0.2cm}
\sourcenote{AStock研究团队测算。评分采用10分制，综合评分为加权得分；标的按板块评级（超配/标配/关注）分组，组内按综合评分降序排列。}
\end{exhibitbox}

\textbf{超配组合（7只）}：功率半导体3只（时代电气、斯达半导、扬杰科技）、电力公用事业2只（长江电力、国电南瑞）、智驾链零部件2只（德赛西威、伯特利）。该组合预计2026年净利润增速中位数约35\%，当前PE（2026E）中位数约22倍，PEG约0.6，具备显著的估值性价比。

\textbf{标配组合（7只）}：半导体设备材料3只（中微公司、北方华创、拓荆科技）、医疗器械2只（迈瑞医疗、联影医疗）、动力电池1只（宁德时代）、创新药1只（恒瑞医药）。该组合预计2026年净利润增速中位数约20\%，当前PE（2026E）中位数约25倍。

\textbf{关注组合（6只）}：包括东微半导、新洁能、华润微、华大智造、药明康德、比亚迪等细分领域龙头。这些标的基本面扎实，景气度向上趋势明确，但短期可能面临估值切换或情绪面扰动，建议逢低布局。

\textbf{组合整体特征}：20只标的合计市值约6.5万亿元，占A股总市值约7\%；2026年预期净利润合计约3,500亿元，平均PE约19倍；2024--2026年净利润复合增速约22\%，显著高于全部A股约8\%的平均水平。从分布来看，制造与科技属性占比约70\%，公用事业与消费属性占比约30\%，整体偏向成长风格。

\section{关键数据速览：板块景气度与估值对比表}

为便于投资者快速把握各板块核心数据，我们整理了重点板块的景气度指标与估值水平对比表。

\begin{exhibitbox}[图表1-4：七大重点板块估值对比表（截至2026年6月22日）]
\scriptsize
\setlength{\tabcolsep}{3pt}
\centering
\vspace{0.3cm}
\begin{tabularx}{\textwidth}{@{}X r r r r r r c @{}}
\toprule
\textbf{板块名称} & \multicolumn{2}{c}{\textbf{PE估值}} & \multicolumn{2}{c}{\textbf{PB估值}} & \textbf{2026E} & \textbf{年初至今} & \textbf{交易} \\
\cmidrule(lr){2-3} \cmidrule(lr){4-5}
& \textbf{PE(TTM)} & \textbf{近5年分位} & \textbf{PB} & \textbf{近5年分位} & \textbf{PEG} & \textbf{涨跌幅} & \textbf{拥挤度} \\
\midrule
功率半导体 & 38.7x & 42.6\% & 4.85x & 35.1\% & 0.87 & +12.4\% & 中 \\
电力/公用事业 & 16.2x & 68.3\% & 1.52x & 71.5\% & 1.94 & +8.7\% & 低 \\
汽车零部件(智驾链) & 32.5x & 31.8\% & 3.46x & 28.4\% & 0.65 & +18.9\% & 高 \\
半导体设备材料 & 47.8x & 55.2\% & 6.12x & 48.7\% & 0.93 & -5.3\% & 中 \\
医疗器械 & 27.4x & 23.5\% & 3.08x & 19.6\% & 0.72 & -2.1\% & 低 \\
动力电池(龙头) & 21.6x & 15.8\% & 2.35x & 12.4\% & 0.58 & -8.6\% & 低 \\
创新药(龙头) & 34.2x & 38.9\% & 3.94x & 42.3\% & 1.15 & +6.2\% & 中 \\
\bottomrule
\end{tabularx}
\vspace{0.2cm}
\sourcenote{Wind、东方财富Choice、AStock研究团队测算。PEG基于2026年一致预期净利润增速计算；分位数统计区间为2021年6月至2026年6月。}
\end{exhibitbox}

\textbf{景气度维度}：功率半导体以28\%的营收增速和3.2个月的库存周期位居首位，显示行业处于高景气且库存健康的理想状态；电力公用事业营收增速仅8\%但利润增速达22\%，体现盈利质量改善逻辑；智驾链零部件营收增速25\%，订单能见度延伸至6个月以上，产业趋势明确；医疗器械营收增速恢复至15\%，库存周期回落至4.5个月，边际改善信号清晰。

\textbf{估值维度}：电力公用事业PE（2026E）仅13倍，处于近五年5\%分位，估值安全边际最高；医疗器械PE约20倍，处于近十年15\%分位，估值修复空间大；功率半导体PE约28倍，处于近五年40\%分位，考虑业绩增速后估值合理；半导体设备PE约35倍，处于近五年60\%分位，估值并不便宜但有国产替代逻辑支撑。

\textbf{持仓维度}：电力公用事业机构持仓比例仅5\%，处于近五年20\%分位，属于典型的低配板块；功率半导体持仓比例约8\%，处于近五年35\%分位，配置水平显著低于景气度排名；医疗器械持仓比例约6\%，处于近十年10\%分位，机构配置意愿处于历史低位；而AI算力板块持仓比例高达25\%，处于近五年95\%分位，交易拥挤度极高。

\textbf{预期差综合评分}：我们将景气度排名与持仓排名的差值定义为“预期差指数”，指数越高代表预期差越大。电力公用事业预期差指数达42，位居第一；功率半导体预期差指数36，排名第二；智驾链零部件预期差指数28，排名第三；而AI算力板块预期差指数为-35，属于预期过度充分的典型。

\section{报告使用指南：研究框架与阅读路径}

本报告采用“总—分—总”的结构安排，按照“框架搭建—板块深度—组合构建—风险提示”的逻辑层层展开，读者可根据自身需求选择不同的阅读路径。

\subsection{研究框架概览}

本报告的核心方法论是\textbf{四维度交叉验证框架}和\textbf{“拥挤度—兑现度”二维矩阵}。第一章系统阐述研究框架与核心结论；第二章阐述研究方法论；第三章进行全市场板块景气度全景扫描；第四至第六章分板块深度剖析，每个板块都按照“景气度验证—预期差分析—产业链图谱—催化与风险”的统一框架展开；第七章警示价值陷阱板块；第八章汇总投资配置建议；第九章提示风险与跟踪指标；第十章为附录。

\begin{exhibitbox}[图表1-5：报告研究框架与章节逻辑图]
\centering
\vspace{0.3cm}
\begin{tikzpicture}[scale=0.85,
  block/.style={rectangle, draw=navy, thick, fill=navy!10, text width=3.2cm, align=center, minimum height=1.0cm, font=\small\bfseries, rounded corners=3pt},
  subblock/.style={rectangle, draw=steelgrey, thick, fill=white, text width=2.8cm, align=center, minimum height=0.8cm, font=\scriptsize, rounded corners=2pt},
  arrow/.style={->, thick, navy}
]

% 第一层：总论
\node[block, fill=navy!20] (ch1) at (0, 6) {第一章 投资概要\\核心结论速览};

% 第二层：方法论
\node[block] (ch2) at (-4, 3.5) {第二章 研究框架\\四维度交叉验证};
\node[block] (ch3) at (4, 3.5) {第三章 全景扫描\\120+板块景气度};

% 第三层：板块深度（2列3行）
\node[subblock] (ch4) at (-5.5, 0.5) {第四章 明星型板块\\功率半导体+电力公用};
\node[subblock] (ch5) at (0, 0.5) {第五章 拐点型板块\\智驾链+半导体设备};
\node[subblock] (ch6) at (5.5, 0.5) {第六章 左侧布局\\医疗器械+动力电池+创新药};

% 第四层：警示与建议
\node[subblock, fill=riskred!10, draw=riskred] (ch7) at (-4, -2.5) {第七章 价值陷阱\\6G+低空经济+卫星互联网};
\node[subblock, fill=nvgreen!15, draw=nvgreen] (ch8) at (0, -2.5) {第八章 配置建议\\组合构建与调仓};
\node[subblock, fill=riskamber!10, draw=riskamber] (ch9) at (4, -2.5) {第九章 风险提示\\跟踪指标体系};

% 第五层：附录
\node[block, fill=panelgrey] (ch10) at (0, -5) {第十章 附录\\数据明细与方法说明};

% 连线
\draw[arrow] (ch1) -- (ch2);
\draw[arrow] (ch1) -- (ch3);
\draw[arrow] (ch2) -- (ch4);
\draw[arrow] (ch2) -- (ch5);
\draw[arrow] (ch2) -- (ch6);
\draw[arrow] (ch3) -- (ch4);
\draw[arrow] (ch3) -- (ch5);
\draw[arrow] (ch3) -- (ch6);
\draw[arrow] (ch4) -- (ch7);
\draw[arrow] (ch5) -- (ch8);
\draw[arrow] (ch6) -- (ch9);
\draw[arrow] (ch7) -- (ch10);
\draw[arrow] (ch8) -- (ch10);
\draw[arrow] (ch9) -- (ch10);

\end{tikzpicture}
\vspace{0.2cm}
\sourcenote{AStock研究团队绘制。}
\end{exhibitbox}

\subsection{推荐阅读路径}

对于投资决策人员，建议采用“\textbf{1\ding{220}8\ding{220}9\ding{220}选读板块}”的快速阅读路径：先读第一章把握核心结论，再读第八章获取具体配置建议，最后读第九章了解风险因素，再根据自身持仓情况选读相关板块的深度章节。

对于行业研究人员，建议采用“\textbf{2\ding{220}3\ding{220}板块深度\ding{220}8}”的深度阅读路径：先读第二章了解整体研究方法与行业比较框架，再读第三章的全景扫描，逐一研读各板块深度章节（第四至第七章），最后读配置建议章节验证判断。

对于量化投资人员，建议重点关注\textbf{第二章的预期差量化指标体系}和\textbf{第八章的组合构建方法}，可直接用于因子开发与策略构建。

\subsection{数据说明与口径}

本报告财务数据主要来源于Wind、Choice、上市公司财报，产业数据来源于行业协会（SEMI、CPCA等）、第三方咨询机构（Yole、Omdia、TrendForce等）及产业链调研，政策数据来源于政府官网与公开信息。估值数据基准日为2026年6月20日。

本报告中的盈利预测为研究团队基于行业景气度判断与公司基本面分析所得，与市场一致预期的差异是我们判断预期差的重要依据。投资者在使用本报告时，应结合自身风险偏好与投资周期，灵活调整配置策略。

\vspace{0.5em}

\begin{riskbox}[本章风险提示]
宏观经济下行超预期导致企业盈利不及预期；产业政策落地进度与力度低于预期；AI技术进展慢于市场预期，影响相关板块景气度；海外制裁升级对半导体等行业供应链造成冲击；市场风格快速切换带来的估值波动风险。
\end{riskbox}
